\documentclass[11pt]{article}

\usepackage[margin=1in]{geometry}
\usepackage{amsmath,amssymb}
\usepackage[T1]{fontenc}
\usepackage{lmodern}
\usepackage{hyperref}
\usepackage{array}
\usepackage{booktabs}

\title{Technical Description of the Fixed-\texorpdfstring{$y$}{y} Quantum Critical-Point Solver}
\author{}
\date{\today}

\begin{document}
\maketitle

\begin{abstract}
This document explains, in plain technical English, what the finite-temperature
fixed-composition solver does in the \texttt{nuclear\_matter\_workflows}
repository, how the code is organized, and where the self-consistent-field
(SCF) iterations appear. The main target is the code in
\texttt{asymmetric\_clausius\_fit/src/asymmetric\_clausius\_fit/finite\_temperature.py},
which computes the critical temperature $T_c$ and critical density $n_c$ for
asymmetric Clausius nuclear matter at fixed proton fraction $y$, using quantum
Fermi statistics. The document also points to the lower-level Fermi-gas module
used by this solver and to the generic SCF implementation on which the
fixed-$y$ version is modeled.
\end{abstract}

\section{Purpose of the Code}

The code solves a finite-temperature critical-point problem for asymmetric
nuclear matter with fixed proton fraction $y$. For a given interaction fit
(parameters $a_n$, $a_{pn}$, $b_n$, $b_{pn}$, and $c$), it searches for the
point $(T,n)$ at which the pressure satisfies the criticality conditions
\begin{align}
\left(\frac{\partial P}{\partial n}\right)_T &= 0, \\
\left(\frac{\partial^2 P}{\partial n^2}\right)_T &= 0.
\end{align}

The output of the solver is therefore the set
\[
(T_c, n_c, P_c),
\]
together with diagnostics such as the residual derivatives, the number of
iterations, and the effective chemical potentials for protons and neutrons.

The finite-temperature part is quantum mechanical because it does not use a
classical Boltzmann gas. Instead, it uses Fermi-Dirac occupation factors and
evaluates the ideal-gas number density and pressure numerically from momentum
integrals.

\section{Relevant Files}

The main files involved in this workflow are listed in
Table~\ref{tab:files}.

\begin{table}[h]
\centering
\begin{tabular}{>{\raggedright\arraybackslash}p{0.42\textwidth} >{\raggedright\arraybackslash}p{0.48\textwidth}}
\toprule
File & Role \\
\midrule
\texttt{asymmetric\_clausius\_fit/src/asymmetric\_clausius\_fit/finite\_temperature.py} & Fixed-$y$ finite-temperature critical-point solver and wrappers. \\
\texttt{quantum\_tc\_nc\_scf/src/quantum\_tc\_nc\_scf/quantum\_fermi.py} & Numerical evaluation of ideal Fermi-gas density and pressure. \\
\texttt{quantum\_tc\_nc\_scf/src/quantum\_tc\_nc\_scf/constants.py} & Physical constants and solver settings. \\
\texttt{quantum\_tc\_nc\_scf/src/quantum\_tc\_nc\_scf/utils/numerics.py} & Low-level numerical utilities such as \texttt{linspace} and Simpson weights. \\
\texttt{quantum\_tc\_nc\_scf/src/quantum\_tc\_nc\_scf/scf.py} & Generic symmetric-matter SCF solver; useful as a reference because the fixed-$y$ solver mirrors its structure. \\
\bottomrule
\end{tabular}
\caption{Core files used by the fixed-$y$ quantum critical-point workflow.}
\label{tab:files}
\end{table}

\section{High-Level Workflow}

At a high level, the solver proceeds as follows:

\begin{enumerate}
\item Read one fitted parameter set from the asymmetric Clausius fit table.
\item Fix the proton fraction $y$.
\item For any trial pair $(T,n)$, construct the corresponding proton and neutron
ideal-gas target densities after excluded-volume corrections.
\item Invert those target densities to effective chemical potentials
$\mu_p^\ast$ and $\mu_n^\ast$ using an inner iterative solver.
\item Use the resulting $\mu_p^\ast$ and $\mu_n^\ast$ to compute the total
pressure $P(T,n;y)$.
\item Approximate $\partial P/\partial n$ and $\partial^2 P/\partial n^2$
numerically.
\item Adjust $(T,n)$ with an outer damped iteration until both criticality
conditions are satisfied.
\end{enumerate}

This means that the code has two nested numerical loops:

\begin{itemize}
\item an \emph{inner} loop that solves for $\mu_p^\ast$ and $\mu_n^\ast$ at
fixed $(T,n)$;
\item an \emph{outer} loop that solves for $(T_c,n_c)$.
\end{itemize}

\section{Data Structures and Helper Functions}

\subsection{\texttt{FixedYCriticalPoint}}

The dataclass \texttt{FixedYCriticalPoint} is a container for the final result.
It stores:
\begin{itemize}
\item the thermodynamic critical point $T_c$, $n_c$, and $P_c$;
\item the residual derivatives at the final point;
\item the convergence score and iteration count;
\item the final effective chemical potentials $\mu_p^\ast$ and $\mu_n^\ast$;
\item the fitted interaction parameters that generated the solution.
\end{itemize}

This object is purely organizational. It does not perform any physics or
numerics by itself.

\subsection{\texttt{\_clamp}}

The helper \texttt{\_clamp(value, lower, upper)} restricts a numerical update to
a safe interval. It is used to prevent the iterative solver from stepping into
unphysical or numerically unstable parameter regions.

\subsection{\texttt{\_attractive\_coefficient}}

The helper
\[
\texttt{\_attractive\_coefficient}(y, a_n, a_{pn})
\]
computes the effective attractive coupling at fixed proton fraction $y$:
\begin{equation}
a_{\text{eff}}(y)
= a_n\left[y^2 + (1-y)^2\right] + 2 a_{pn} y (1-y).
\end{equation}

This is the coefficient that enters the interaction pressure term.

\subsection{\texttt{\_species\_data}}

The helper \texttt{\_species\_data(n, y, b\_n, b\_{pn})} converts the total
baryon density $n$ and proton fraction $y$ into species-specific quantities:
\begin{align}
n_p &= yn, \\
n_n &= (1-y)n.
\end{align}

It then builds excluded-volume combinations
\begin{align}
x_p &= b_n n_p + b_{pn} n_n, \\
x_n &= b_{pn} n_p + b_n n_n,
\end{align}
followed by free-volume factors
\begin{align}
f_p &= 1 - x_p, \\
f_n &= 1 - x_n.
\end{align}

If either $f_p \le 0$ or $f_n \le 0$, the state is rejected as unphysical.
Otherwise, the code defines target ideal-gas densities
\begin{align}
n_{p,\mathrm{id}}^{\mathrm{target}} &= \frac{n_p}{f_p}, \\
n_{n,\mathrm{id}}^{\mathrm{target}} &= \frac{n_n}{f_n}.
\end{align}

These target densities are the quantities that the inner SCF loop must match.

\section{The Quantum Fermi-Gas Module}

The file \texttt{quantum\_fermi.py} defines the class
\texttt{QuantumFermiGas}. This class is responsible for evaluating the ideal
finite-temperature number density and pressure of a relativistic Fermi gas.

\subsection{Momentum Grid and Numerical Integration}

When a \texttt{QuantumFermiGas} object is created, it precomputes:
\begin{itemize}
\item a momentum grid $k$;
\item Simpson integration weights;
\item the single-particle energy
\[
E(k) = \sqrt{(\hbar c\, k)^2 + m^2};
\]
\item the momentum-dependent factors needed for density and pressure
integrals.
\end{itemize}

The integration grid is produced by \texttt{simpson\_weights} from
\texttt{utils/numerics.py}. The reason for this precomputation is efficiency:
the same quadrature grid is reused many times during the iterative solution.

\subsection{Fermi-Dirac Occupation}

The occupation factor is
\begin{equation}
f(E;\mu,T) = \frac{1}{e^{(E-\mu)/T} + 1},
\end{equation}
with overflow protection for very large positive or negative arguments.

\subsection{Ideal Density and Pressure}

The code numerically evaluates
\begin{align}
n_{\mathrm{id}}(T,\mu) &=
\frac{d}{2\pi^2} \int_0^{k_{\max}} k^2 f(E;\mu,T)\, dk, \\
p_{\mathrm{id}}(T,\mu) &=
\frac{d}{6\pi^2} \int_0^{k_{\max}}
\frac{(\hbar c)^2 k^4}{E(k)} f(E;\mu,T)\, dk,
\end{align}
where $d$ is the degeneracy factor.

In the fixed-$y$ solver, the settings are modified so that each species uses
$d=2$ rather than $d=4$. This is because the proton and neutron sectors are
treated separately.

\subsection{Initial Guess for \texorpdfstring{$\mu^\ast$}{mu*}}

The method \texttt{mu\_seed\_from\_n\_id} provides a zero-temperature estimate
for the effective chemical potential. It computes a Fermi momentum from the
target ideal density and then returns the corresponding relativistic Fermi
energy:
\begin{align}
k_F &= \left(\frac{6\pi^2 n_{\mathrm{id}}}{d}\right)^{1/3}, \\
\mu_{\text{seed}} &= \sqrt{(\hbar c\, k_F)^2 + m^2}.
\end{align}

This serves as the starting point for the inner iterative inversion.

\section{Pressure Evaluator for Fixed Composition}

The class \texttt{FixedYAsymmetricClausiusPressureEvaluator} is the central
object in the fixed-$y$ workflow.

\subsection{Initialization}

The constructor receives the interaction parameters
\[
a_n,\quad a_{pn},\quad b_n,\quad b_{pn},\quad c,
\]
the fixed proton fraction $y$, and a settings object. It creates two
\texttt{QuantumFermiGas} instances:
\begin{itemize}
\item one for protons;
\item one for neutrons.
\end{itemize}

It also creates a cache keyed by $(T,n)$ so that repeated calls to the same
thermodynamic point do not recompute the pressure.

\subsection{Target Ideal Densities}

The method \texttt{target\_n\_id(n)} calls \texttt{\_species\_data} and returns
the pair
\[
\left(n_{p,\mathrm{id}}^{\mathrm{target}},\;
n_{n,\mathrm{id}}^{\mathrm{target}}\right).
\]

This is the information needed to perform the inversion from density to
effective chemical potential.

\section{Where the SCF Appears}

This is the most important conceptual point. The code contains \emph{two}
different SCF layers.

\subsection{Inner SCF: Solving for \texorpdfstring{$\mu_p^\ast$}{mu*p} and \texorpdfstring{$\mu_n^\ast$}{mu*n}}

The inner SCF appears in the method
\texttt{\_solve\_mu\_star\_species}. For a fixed temperature $T$ and a fixed
target ideal density $n_{\mathrm{id}}^{\mathrm{target}}$, the code solves
\begin{equation}
F(\mu^\ast) =
n_{\mathrm{id}}(T,\mu^\ast) - n_{\mathrm{id}}^{\mathrm{target}} = 0.
\end{equation}

The loop starts from the zero-temperature seed and then repeatedly:
\begin{enumerate}
\item evaluates the residual
\[
r = n_{\mathrm{id}}(T,\mu) - n_{\mathrm{id}}^{\mathrm{target}};
\]
\item estimates the derivative numerically,
\begin{equation}
\chi \approx
\frac{n_{\mathrm{id}}(T,\mu+\delta\mu) -
n_{\mathrm{id}}(T,\mu-\delta\mu)}
{(\mu+\delta\mu)-(\mu-\delta\mu)};
\end{equation}
\item computes the update
\begin{equation}
\Delta\mu = \frac{r}{\chi};
\end{equation}
\item applies damping,
\begin{equation}
\mu_{\mathrm{new}} = \mu - \lambda_\mu \Delta\mu,
\end{equation}
where $\lambda_\mu = \texttt{mu\_scf\_damping}$;
\item clamps the result to a safe interval and repeats.
\end{enumerate}

Strictly speaking, this is a damped Newton step with a numerical derivative,
embedded in a self-consistent loop. For that reason, the code calls it SCF,
but mathematically it is closer to a damped Newton-Raphson inversion than to a
pure Picard fixed-point iteration.

\subsection{Outer SCF: Solving for \texorpdfstring{$(T_c,n_c)$}{(Tc,nc)}}

The outer SCF appears in the function
\texttt{solve\_critical\_point\_fixed\_y\_scf}. This solver seeks the pair
$(T,n)$ satisfying
\begin{align}
G_1(T,n) &= \left(\frac{\partial P}{\partial n}\right)_T = 0, \\
G_2(T,n) &= \left(\frac{\partial^2 P}{\partial n^2}\right)_T = 0.
\end{align}

The procedure is:
\begin{enumerate}
\item obtain a coarse initial seed $(T,n)$ from a grid scan;
\item evaluate $G_1$ and $G_2$ at the current point;
\item estimate the sensitivity of $G_1$ to $T$ using finite differences,
\begin{equation}
\frac{\partial G_1}{\partial T}
\approx
\frac{G_1(T+h_T,n)-G_1(T-h_T,n)}{2h_T};
\end{equation}
\item estimate the sensitivity of $G_2$ to $n$ using finite differences,
\begin{equation}
\frac{\partial G_2}{\partial n}
\approx
\frac{G_2(T,n+h_n)-G_2(T,n-h_n)}{2h_n};
\end{equation}
\item form damped correction steps
\begin{align}
\Delta T &= -\lambda_T \frac{G_1}{\partial G_1/\partial T}, \\
\Delta n &= -\lambda_n \frac{G_2}{\partial G_2/\partial n};
\end{align}
\item clamp the step sizes to avoid excessive moves;
\item run a backtracking line search and accept the trial point only if it
reduces the normalized residual score.
\end{enumerate}

This is again SCF in the broad sense of a self-consistent outer iteration, but
its update formula is Newton-like because it uses local slope information.

\section{How the Pressure is Constructed}

Once the inner loop has produced $\mu_p^\ast$ and $\mu_n^\ast$, the method
\texttt{pressure(T,n)} evaluates:
\begin{enumerate}
\item proton ideal pressure $p_p(T,\mu_p^\ast)$;
\item neutron ideal pressure $p_n(T,\mu_n^\ast)$;
\item interaction contribution
\begin{equation}
p_{\mathrm{int}}(n,y) =
-\frac{a_{\mathrm{eff}}(y)\, n^2}{(1 + cn)^2}.
\end{equation}
\end{enumerate}

The total pressure is then
\begin{equation}
P(T,n;y) = p_p(T,\mu_p^\ast) + p_n(T,\mu_n^\ast) + p_{\mathrm{int}}(n,y).
\end{equation}

This is the function whose density derivatives define the critical point.

\section{Numerical Derivatives Used by the Critical Solver}

The method \texttt{critical\_equations(T,n)} computes the first and second
density derivatives numerically using central differences:
\begin{align}
\frac{\partial P}{\partial n}
&\approx \frac{P(T,n+h_n)-P(T,n-h_n)}{2h_n}, \\
\frac{\partial^2 P}{\partial n^2}
&\approx \frac{P(T,n+h_n)-2P(T,n)+P(T,n-h_n)}{h_n^2}.
\end{align}

The code checks that all three pressure evaluations exist and remain inside the
allowed density domain. If any evaluation fails, the solver returns
\texttt{None} for that point.

\section{Seed Search Before the Outer Iteration}

The function \texttt{find\_seed\_point\_fixed\_y} performs a coarse search over
a rectangular grid in temperature and density. For each valid point it computes
$G_1$ and $G_2$, then defines a normalized score
\begin{equation}
S(T,n) =
\left(\frac{G_1}{s_1}\right)^2 +
\left(\frac{G_2}{s_2}\right)^2,
\end{equation}
where $s_1$ and $s_2$ are the largest absolute values found on the coarse grid.

The best seed is the point with the smallest score. This is a practical way to
start the outer solver close enough to the critical point for the damped
iteration to converge.

\section{Control Flow of the Main Public Functions}

\subsection{\texttt{compute\_fixed\_y\_critical\_point\_from\_fit\_row}}

This function is the main single-case wrapper. It:
\begin{enumerate}
\item receives one row from a fit table;
\item creates a \texttt{FixedYAsymmetricClausiusPressureEvaluator};
\item calls \texttt{solve\_critical\_point\_fixed\_y\_scf};
\item packages the returned values into a \texttt{FixedYCriticalPoint}.
\end{enumerate}

This is the function that a user should think of as ``solve one critical point
for one fitted model at one fixed composition''.

\subsection{\texttt{compute\_fixed\_y\_family\_from\_fit\_df}}

This function loops over many fit rows and many proton fractions $y$. It uses a
continuation strategy:
\begin{itemize}
\item if the previous $y$ value converged, the solver reuses
$(T_c,n_c)$ as the next seed;
\item if one point fails, the continuation seed is reset.
\end{itemize}

This helps smooth parameter sweeps across composition.

\section{Settings That Control the Solver}

The file \texttt{constants.py} defines the numerical settings. The most
important ones are:
\begin{itemize}
\item \texttt{mu\_scf\_tol}, \texttt{mu\_scf\_max\_iter},
\texttt{mu\_scf\_damping}, and \texttt{mu\_delta} for the inner
$\mu^\ast$ inversion;
\item \texttt{h\_n} and \texttt{h\_t} for finite-difference derivatives;
\item \texttt{outer\_tol\_dPdn} and \texttt{outer\_tol\_d2Pdn2} for the final
criticality tolerances;
\item \texttt{outer\_damping\_t}, \texttt{outer\_damping\_n},
\texttt{outer\_max\_step\_t}, \texttt{outer\_max\_step\_n}, and
\texttt{outer\_min\_lambda} for the outer update and line search;
\item \texttt{coarse\_t\_count} and \texttt{coarse\_n\_count} for the initial
seed grid.
\end{itemize}

The settings file therefore encodes both the physical scale of the search
window and the numerical stability rules of the algorithm.

\section{Why the Code Returns \texttt{None}}

The solver uses \texttt{None} to signal failure rather than forcing a result.
Typical failure modes are:
\begin{itemize}
\item an excluded-volume factor becomes non-positive;
\item a finite-difference slope is too small to be reliable;
\item the inner inversion for $\mu^\ast$ does not converge;
\item the line search cannot find an improving trial point;
\item the requested point lies outside the allowed temperature or density range.
\end{itemize}

This makes the code conservative: it prefers returning no solution to returning
a numerically dubious one.

\section{Relationship to the Generic Quantum SCF Solver}

The file \texttt{quantum\_tc\_nc\_scf/src/quantum\_tc\_nc\_scf/scf.py}
contains a more generic version of the same strategy for symmetric models. The
fixed-$y$ solver in \texttt{finite\_temperature.py} follows the same overall
design:
\begin{itemize}
\item an evaluator object that computes pressure;
\item an inner inversion from density to $\mu^\ast$;
\item a coarse seed search;
\item an outer damped iteration for the critical conditions.
\end{itemize}

The asymmetric fixed-$y$ solver extends this pattern by splitting the matter
into proton and neutron sectors and by introducing the composition-dependent
interaction coefficient.

\section{Summary}

The fixed-$y$ finite-temperature solver computes critical points for asymmetric
Clausius nuclear matter using quantum Fermi statistics and nested iterative
numerics.

Its essential logic is:
\begin{enumerate}
\item map the physical density $n$ into proton and neutron target ideal
densities through excluded-volume relations;
\item invert each target density to an effective chemical potential using an
inner damped Newton-like SCF loop;
\item build the total pressure from proton, neutron, and interaction terms;
\item compute the first and second density derivatives numerically;
\item adjust $(T,n)$ with an outer damped Newton-like SCF loop until the
criticality conditions are satisfied.
\end{enumerate}

Therefore, the answer to ``where is the SCF in the code?'' is:
\begin{itemize}
\item \textbf{inner SCF:} in
\texttt{\_solve\_mu\_star\_species}, where $\mu_p^\ast$ and $\mu_n^\ast$ are
found from target densities;
\item \textbf{outer SCF:} in
\texttt{solve\_critical\_point\_fixed\_y\_scf}, where $T_c$ and $n_c$ are found
from the vanishing of $\partial P/\partial n$ and $\partial^2 P/\partial n^2$.
\end{itemize}

These two layers are the conceptual heart of the code.

\end{document}
